% Numeric bibliography in first-citation order.
\begin{thebibliography}{99}
\setlength{\itemsep}{0.35em}

\def\currentbiburl{https://doi.org/10.1186/s12874-023-02078-1}
\bibitem{ref1} Y. Huang, J. Li, M. Li, and R. R. Aparasu, ``Application of machine learning in predicting survival outcomes involving real-world data: a scoping review,'' BMC Medical Research Methodology, vol. 23, no. 1, p. 268, Nov. 2023, doi: \href{https://doi.org/10.1186/s12874-023-02078-1}{10.1186/s12874-023-02078-1}

\def\currentbiburl{https://doi.org/10.1145/3214306}
\bibitem{ref3} P. Wang, Y. Li, and C. K. Reddy, ``Machine learning for survival analysis,'' ACM Computing Surveys, vol. 51, no. 6, pp. 1--36, Feb. 2019, doi: \href{https://doi.org/10.1145/3214306}{10.1145/3214306}

\def\currentbiburl{https://doi.org/10.4097/kja.d.18.00067}
\bibitem{ref14} J. In and D. K. Lee, ``Survival analysis: Part I: analysis of time-to-event,'' Korean Journal of Anesthesiology, vol. 71, no. 3, pp. 182--191, May 2018, doi: \href{https://doi.org/10.4097/kja.d.18.00067}{10.4097/kja.d.18.00067}

\def\currentbiburl{https://doi.org/10.1155/2022/1176060}
\bibitem{ref5} G. Kantidakis, A.-D. Hazewinkel, and M. Fiocco, ``Neural Networks for Survival Prediction in Medicine Using Prognostic Factors: A Review and Critical Appraisal,'' Computational and Mathematical Methods in Medicine, vol. 2022, pp. 1--18, Sep. 2022, doi: \href{https://doi.org/10.1155/2022/1176060}{10.1155/2022/1176060}

\def\currentbiburl{https://doi.org/10.1186/s12874-018-0482-1}
\bibitem{ref4} J. L. Katzman, U. Shaham, A. Cloninger, J. Bates, T. Jiang, and Y. Kluger, ``DeepSurv: personalized treatment recommender system using a Cox proportional hazards deep neural network,'' BMC Medical Research Methodology, vol. 18, no. 1, p. 24, Feb. 2018, doi: \href{https://doi.org/10.1186/s12874-018-0482-1}{10.1186/s12874-018-0482-1}

\def\currentbiburl{https://doi.org/10.1016/j.jbi.2021.103959}
\bibitem{ref2} F. Xie et al., ``AutoScore-Survival: Developing interpretable machine learning-based time-to-event scores with right-censored survival data,'' Journal of Biomedical Informatics, vol. 125, p. 103959, Nov. 2021, doi: \href{https://doi.org/10.1016/j.jbi.2021.103959}{10.1016/j.jbi.2021.103959}

\def\currentbiburl{https://doi.org/10.1007/s10985-021-09532-6}
\bibitem{ref21} H. Kvamme and Ø. Borgan, ``Continuous and discrete-time survival prediction with neural networks,'' Lifetime Data Analysis, vol. 27, no. 4, pp. 710--736, Oct. 2021, doi: \href{https://doi.org/10.1007/s10985-021-09532-6}{10.1007/s10985-021-09532-6}

\def\currentbiburl{https://doi.org/10.1214/20-aoas1325}
\bibitem{ref6} A. Nikooienejad, W. Wang, and V. E. Johnson, ``Bayesian variable selection for survival data using inverse moment priors,'' The Annals of Applied Statistics, vol. 14, no. 2, pp. 809--828, Jun. 2020, doi: \href{https://doi.org/10.1214/20-aoas1325}{10.1214/20-aoas1325}

\def\currentbiburl{https://doi.org/10.1002/sim.8933}
\bibitem{ref22} D. Alvares, E. Lázaro, V. Gómez\textendash{}Rubio, and C. Armero, ``Bayesian survival analysis with BUGS,'' Statistics in Medicine, vol. 40, no. 12, pp. 2975--3020, Mar. 2021, doi: \href{https://doi.org/10.1002/sim.8933}{10.1002/sim.8933}

\def\currentbiburl{https://doi.org/10.1002/cnr2.1176}
\bibitem{ref15} M. Perera and A. K. Dwivedi, ``Statistical issues and methods in designing and analyzing survival studies,'' Cancer Reports, vol. 3, no. 4, p. e1176, 2020, doi: \href{https://doi.org/10.1002/cnr2.1176}{10.1002/cnr2.1176}

\def\currentbiburl{https://doi.org/10.1186/s12874-019-0749-1}
\bibitem{ref17} K. Jachno, S. Heritier, and R. Wolfe, ``Are non-constant rates and non-proportional treatment effects accounted for in the design and analysis of randomised controlled trials? A review of current practice,'' BMC Medical Research Methodology, vol. 19, no. 1, p. 103, May 2019, doi: \href{https://doi.org/10.1186/s12874-019-0749-1}{10.1186/s12874-019-0749-1}

\def\currentbiburl{https://doi.org/10.7326/m22-0844}
\bibitem{ref84} D. J. McLernon et al., ``Assessing performance and clinical usefulness in prediction models with survival outcomes: Practical guidance for COX Proportional Hazards Models,'' Annals of Internal Medicine, vol. 176, no. 1, pp. 105--114, 2023, doi: \href{https://doi.org/10.7326/m22-0844}{10.7326/m22-0844}

\def\currentbiburl{https://doi.org/10.3348/kjr.2021.0223}
\bibitem{ref26} S. Y. Park, J. E. Park, H. Kim, and S. H. Park, ``Review of Statistical Methods for Evaluating the Performance of Survival or Other Time-to-Event Prediction Models (from Conventional to Deep Learning Approaches),'' Korean Journal of Radiology, vol. 22, no. 10, p. 1697, Jan. 2021, doi: \href{https://doi.org/10.3348/kjr.2021.0223}{10.3348/kjr.2021.0223}

\def\currentbiburl{https://doi.org/10.1016/j.jbi.2020.103496}
\bibitem{ref27} E. Longato, M. Vettoretti, and B. Di Camillo, ``A practical perspective on the concordance index for the evaluation and selection of prognostic time-to-event models,'' Journal of Biomedical Informatics, vol. 108, p. 103496, Jul. 2020, doi: \href{https://doi.org/10.1016/j.jbi.2020.103496}{10.1016/j.jbi.2020.103496}

\def\currentbiburl{https://doi.org/10.1002/sim.8570}
\bibitem{ref59} P. C. Austin, F. E. Harrell, and D. van Klaveren, ``Graphical calibration curves and the integrated calibration index (ICI) for survival models,'' Statistics in Medicine, vol. 39, no. 21, pp. 2714--2742, Jun. 2020, doi: \href{https://doi.org/10.1002/sim.8570}{10.1002/sim.8570}

\def\currentbiburl{https://doi.org/10.1093/gigascience/giac071}
\bibitem{ref33} Y. Zhang, G. Wong, G. Mann, S. Muller, and J. Y. H. Yang, ``SurvBenchmark: comprehensive benchmarking study of survival analysis methods using both omics data and clinical data,'' GigaScience, vol. 11, p. giac071, Jul. 2022, doi: \href{https://doi.org/10.1093/gigascience/giac071}{10.1093/gigascience/giac071}

\def\currentbiburl{https://doi.org/10.1186/s41512-022-00124-y}
\bibitem{ref44} H. Smith, M. Sweeting, T. Morris, and M. J. Crowther, ``A scoping methodological review of simulation studies comparing statistical and machine learning approaches to risk prediction for time-to-event data,'' Diagnostic and Prognostic Research, vol. 6, no. 1, p. 10, Jun. 2022, doi: \href{https://doi.org/10.1186/s41512-022-00124-y}{10.1186/s41512-022-00124-y}

\def\currentbiburl{https://doi.org/10.1016/j.eswa.2016.07.018}
\bibitem{ref19} M. Schmid, M. N. Wright, and A. Ziegler, ``On the use of Harrell's C for clinical risk prediction via random survival forests,'' Expert Systems With Applications, vol. 63, pp. 450--459, Jul. 2016, doi: \href{https://doi.org/10.1016/j.eswa.2016.07.018}{10.1016/j.eswa.2016.07.018}

\def\currentbiburl{https://doi.org/10.1002/bimj.202100380}
\bibitem{ref37} V. Baralou, N. Kalpourtzi, and G. Touloumi, ``Individual risk prediction: Comparing random forests with Cox proportional\textendash{}hazards model by a simulation study,'' Biometrical Journal, vol. 65, no. 6, p. e2100380, 2023, doi: \href{https://doi.org/10.1002/bimj.202100380}{10.1002/bimj.202100380}

\def\currentbiburl{https://doi.org/10.1002/bimj.201900211}
\bibitem{ref23} E. Lázaro, C. Armero, and D. Alvares, ``Bayesian regularization for flexible baseline hazard functions in Cox survival models,'' Biometrical Journal, vol. 63, no. 1, pp. 7--26, 2021, doi: \href{https://doi.org/10.1002/bimj.201900211}{10.1002/bimj.201900211}

\def\currentbiburl{https://doi.org/10.1002/sim.10160}
\bibitem{ref90} D. Alvares, J. van Niekerk, E. T. Krainski, H. Rue, and D. Rustand, ``Bayesian survival analysis with INLA,'' Statistics in Medicine, vol. 43, no. 20, pp. 3975--4010, Jun. 2024, doi: \href{https://doi.org/10.1002/sim.10160}{10.1002/sim.10160}

\def\currentbiburl{https://doi.org/10.1093/bioinformatics/btx300}
\bibitem{ref25} Z. Tang, Y. Shen, X. Zhang, and N. Yi, ``The spike-and-slab lasso Cox model for survival prediction and associated genes detection,'' Bioinformatics, vol. 33, no. 18, pp. 2799--2807, May 2017, doi: \href{https://doi.org/10.1093/bioinformatics/btx300}{10.1093/bioinformatics/btx300}

\def\currentbiburl{https://doi.org/10.1186/s12859-022-04770-3}
\bibitem{ref42} S. Samorodnitsky, K. A. Hoadley, and E. F. Lock, ``A hierarchical spike-and-slab model for pan-cancer survival using pan-omic data,'' BMC Bioinformatics, vol. 23, no. 1, p. 235, Jun. 2022, doi: \href{https://doi.org/10.1186/s12859-022-04770-3}{10.1186/s12859-022-04770-3}

\def\currentbiburl{https://doi.org/10.1002/pst.2256}
\bibitem{ref38} Y. Qing, P. F. Thall, and Y. Yuan, ``A Bayesian piecewise exponential phase II design for monitoring a time\textendash{}to\textendash{}event endpoint,'' Pharmaceutical Statistics, vol. 22, no. 1, pp. 34--44, 2023, doi: \href{https://doi.org/10.1002/pst.2256}{10.1002/pst.2256}

\def\currentbiburl{https://doi.org/10.1002/jrsm.1722}
\bibitem{ref92} S. C. Freeman, A. J. Sutton, N. J. Cooper, A. Gasparini, M. J. Crowther, and N. Hawkins, ``Bayesian pairwise meta\textendash{}analysis of time\textendash{}to\textendash{}event outcomes in the presence of non\textendash{}proportional hazards: A simulation study of flexible parametric, piecewise exponential and fractional polynomial models,'' Research Synthesis Methods, vol. 15, no. 5, pp. 780--801, May 2024, doi: \href{https://doi.org/10.1002/jrsm.1722}{10.1002/jrsm.1722}

\def\currentbiburl{https://doi.org/10.1016/j.knosys.2022.110234}
\bibitem{ref30} M. Krzyziński, M. Spytek, H. Baniecki, and P. Biecek, ``SurvSHAP(t): Time-dependent explanations of machine learning survival models,'' Knowledge-Based Systems, vol. 262, p. 110234, 2023, doi: \href{https://doi.org/10.1016/j.knosys.2022.110234}{10.1016/j.knosys.2022.110234}

\def\currentbiburl{https://doi.org/10.1002/sim.10057}
\bibitem{ref60} K. Suresh, C. Görg, and D. Ghosh, ``Model\textendash{}agnostic explanations for survival prediction models,'' Statistics in Medicine, vol. 43, no. 11, pp. 2161--2182, Mar. 2024, doi: \href{https://doi.org/10.1002/sim.10057}{10.1002/sim.10057}

\def\currentbiburl{https://doi.org/10.1038/s43586-020-00001-2}
\bibitem{ref40} R. Van De Schoot et al., ``Bayesian statistics and modelling,'' Nature Reviews Methods Primers, vol. 1, no. 1, p. 1, Jan. 2021, doi: \href{https://doi.org/10.1038/s43586-020-00001-2}{10.1038/s43586-020-00001-2}

\def\currentbiburl{https://doi.org/10.1093/biomet/asq017}
\bibitem{carvalho2010}
C. M. Carvalho, N. G. Polson, and J. G. Scott, ``The horseshoe estimator for sparse signals,'' \emph{Biometrika}, vol. 97, no. 2, pp. 465--480, 2010, doi:~\url{https://doi.org/10.1093/biomet/asq017}.

\def\currentbiburl{https://doi.org/10.1214/17-EJS1337SI}
\bibitem{piironen2017}
J. Piironen and A. Vehtari, ``Sparsity Information and Regularization in the Horseshoe and Other Shrinkage Priors,'' \emph{Electronic Journal of Statistics}, vol. 11, no. 2, pp. 5018--5051, 2017, doi:~\href{https://doi.org/10.1214/17-EJS1337SI}{\nolinkurl{10.1214/17-EJS1337SI}}.

\def\currentbiburl{https://doi.org/10.1080/02664763.2021.1893285}
\bibitem{ref43} R. Wu, M. Ahn, and H. Yang, ``Spike-and-slab type variable selection in the Cox proportional hazards model for high-dimensional features,'' Journal of Applied Statistics, vol. 49, no. 9, pp. 2189--2207, 2022, doi: \href{https://doi.org/10.1080/02664763.2021.1893285}{10.1080/02664763.2021.1893285}

\def\currentbiburl{https://doi.org/10.1037/met0000275}
\bibitem{ref53} D. J. Schad, M. Betancourt, and S. Vasishth, ``Toward a principled Bayesian workflow in cognitive science,'' Psychological Methods, vol. 26, no. 1, pp. 103--126, 2021, doi: \href{https://doi.org/10.1037/met0000275}{10.1037/met0000275}

\def\currentbiburl{https://doi.org/10.1214/20-ba1221}
\bibitem{ref39} A. Vehtari, A. Gelman, D. Simpson, B. Carpenter, and P.-C. Bürkner, ``Rank-Normalization, Folding, and Localization: An Improved $\widehat{R}$ for Assessing Convergence of MCMC (with Discussion),'' Bayesian Analysis, vol. 16, no. 2, pp. 667--718, Jun. 2021, doi: \href{https://doi.org/10.1214/20-ba1221}{10.1214/20-ba1221}

\def\currentbiburl{https://doi.org/10.1200/JCO.1994.12.10.2086}
\bibitem{schumacher1994}
M. Schumacher~\emph{et al.}, ``Randomized 2 $\times$ 2 trial evaluating hormonal treatment and the duration of chemotherapy in node-positive breast cancer patients. German Breast Cancer Study Group,'' \emph{Journal of Clinical Oncology}, vol. 12, no. 10, pp. 2086--2093, 1994, doi:~\url{https://doi.org/10.1200/JCO.1994.12.10.2086}.

\def\currentbiburl{https://www.jmlr.org/papers/v21/20-729.html}
\bibitem{polsterl2020}
S. P{\"o}lsterl, ``scikit-survival: a library for time-to-event analysis built on top of scikit-learn,'' \emph{Journal of Machine Learning Research}, vol. 21, no. 212, pp. 1--6, 2020.~\url{https://www.jmlr.org/papers/v21/20-729.html}.

\def\currentbiburl{https://doi.org/10.1038/nature10983}
\bibitem{curtis2012}
C. Curtis~\emph{et al.}, ``The genomic and transcriptomic architecture of 2,000 breast tumours reveals novel subgroups,'' \emph{Nature}, vol. 486, no. 7403, pp. 346--352, 2012, doi:~\url{https://doi.org/10.1038/nature10983}.

\def\currentbiburl{https://doi.org/10.1038/ncomms11479}
\bibitem{pereira2016}
B. Pereira~\emph{et al.}, ``The somatic mutation profiles of 2,433 breast cancers refine their genomic and transcriptomic landscapes,'' \emph{Nature Communications}, vol. 7, no. 1, p. 11479, 2016, doi:~\url{https://doi.org/10.1038/ncomms11479}.

\def\currentbiburl{https://doi.org/10.1002/9780470258019}
\bibitem{hosmer2008}
D. W. Hosmer, S. Lemeshow, and S. May, \emph{Applied Survival Analysis: Regression Modeling of Time-to-Event Data}, 2nd ed. Hoboken, NJ: Wiley, 2008, doi:~\url{https://doi.org/10.1002/9780470258019}.

\def\currentbiburl{https://doi.org/10.1093/bioinformatics/btaa046}
\bibitem{ref32} S. Parvandeh, H.-W. Yeh, M. P. Paulus, and B. A. McKinney, ``Consensus features nested cross-validation,'' Bioinformatics, vol. 36, no. 10, pp. 3093--3098, Jan. 2020, doi: \href{https://doi.org/10.1093/bioinformatics/btaa046}{10.1093/bioinformatics/btaa046}

\def\currentbiburl{https://doi.org/10.1186/s12874-021-01469-6}
\bibitem{ref34} C. L. Andaur Navarro et al., ``Completeness of reporting of clinical prediction models developed using supervised machine learning: a systematic review,'' BMC Medical Research Methodology, vol. 22, no. 1, p. 12, Jan. 2022, doi: \href{https://doi.org/10.1186/s12874-021-01469-6}{10.1186/s12874-021-01469-6}

\def\currentbiburl{https://doi.org/10.1155/2021/9307475}
\bibitem{ref45} J. Turkson, F. Ayiah-Mensah, and V. Nimoh, ``Handling Censoring and Censored Data in Survival Analysis: A Standalone Systematic Literature review,'' International Journal of Mathematics and Mathematical Sciences, vol. 2021, no. 1, Jan. 2021, doi: \href{https://doi.org/10.1155/2021/9307475}{10.1155/2021/9307475}

\def\currentbiburl{https://doi.org/10.1002/sim.9216}
\bibitem{ref47} M. M. Nuño and D. L. Gillen, ``Censoring\textendash{}robust time\textendash{}dependent receiver operating characteristic curve estimators,'' Statistics in Medicine, vol. 40, no. 30, pp. 6885--6899, Oct. 2021, doi: \href{https://doi.org/10.1002/sim.9216}{10.1002/sim.9216}

\def\currentbiburl{https://doi.org/10.1093/biomet/69.1.239}
\bibitem{schoenfeld1982}
D. Schoenfeld, ``Partial residuals for the proportional hazards regression model,'' \emph{Biometrika}, vol. 69, no. 1, pp. 239--241, 1982, doi:~\url{https://doi.org/10.1093/biomet/69.1.239}.

\def\currentbiburl{https://doi.org/10.2307/4615733}
\bibitem{holm1979}
S. Holm, ``A simple sequentially rejective multiple test procedure,'' \emph{Scandinavian Journal of Statistics}, vol. 6, no. 2, pp. 65--70, 1979, doi:~\url{https://doi.org/10.2307/4615733}.

\def\currentbiburl{https://doi.org/10.22283/qbs.2017.36.2.85}
\bibitem{ref51} H. Wang and G. Li, ``A selective review on random survival forests for high dimensional data,'' Quantitative Bio-Science, vol. 36, no. 2, pp. 85--96, Nov. 2017, doi: \href{https://doi.org/10.22283/qbs.2017.36.2.85}{10.22283/qbs.2017.36.2.85}

\def\currentbiburl{https://doi.org/10.1080/01621459.1981.10477634}
\bibitem{laird1981}
N. Laird and D. Olivier, ``Covariance analysis of censored survival data using log-linear analysis techniques,'' \emph{Journal of the American Statistical Association}, vol. 76, no. 374, pp. 231--240, 1981, doi:~\url{https://doi.org/10.1080/01621459.1981.10477634}.

\def\currentbiburl{https://doi.org/10.1214/aos/1176345693}
\bibitem{friedman1982}
M. Friedman, ``Piecewise exponential models for survival data with covariates,'' \emph{The Annals of Statistics}, vol. 10, no. 1, pp. 101--113, 1982, doi:~\url{https://doi.org/10.1214/aos/1176345693}.

\def\currentbiburl{https://doi.org/10.1080/00031305.2020.1865198}
\bibitem{ref55} S. Thomas and W. Tu, ``Learning Hamiltonian Monte Carlo in R,'' The American Statistician, vol. 75, no. 4, pp. 403--413, 2021, doi: \href{https://doi.org/10.1080/00031305.2020.1865198}{10.1080/00031305.2020.1865198}

\def\currentbiburl{https://doi.org/10.1146/annurev-statistics-040220-091727}
\bibitem{ref54} L. F. South, M. Riabiz, O. Teymur, and C. J. Oates, ``Postprocessing of MCMC,'' Annual Review of Statistics and Its Application, vol. 9, pp. 529--555, Mar. 2022, doi: \href{https://doi.org/10.1146/annurev-statistics-040220-091727}{10.1146/annurev-statistics-040220-091727}

\def\currentbiburl{https://doi.org/10.1023/A:1024068626366}
\bibitem{nadeau2003}
C. Nadeau and Y. Bengio, ``Inference for the generalization error,'' \emph{Machine Learning}, vol. 52, no. 3, pp. 239--281, 2003, doi:~\url{https://doi.org/10.1023/A:1024068626366}.

\def\currentbiburl{https://doi.org/10.1038/s41598-023-35795-0}
\bibitem{ref28} R. O. Alabi, M. Elmusrati, I. Leivo, A. Almangush, and A. A. Mäkitie, ``Machine learning explainability in nasopharyngeal cancer survival using LIME and SHAP,'' Scientific Reports, vol. 13, no. 1, p. 8984, Jun. 2023, doi: \href{https://doi.org/10.1038/s41598-023-35795-0}{10.1038/s41598-023-35795-0}

\def\currentbiburl{https://doi.org/10.1111/biom.13124}
\bibitem{ref56} M. Wiesenfarth and S. Calderazzo, ``Quantification of prior impact in terms of effective current sample size,'' Biometrics, vol. 76, no. 1, pp. 326--336, 2020, doi: \href{https://doi.org/10.1111/biom.13124}{10.1111/biom.13124}

\def\currentbiburl{https://papers.nips.cc/paper/2017/hash/8a20a8621978632d76c43dfd28b67767-Abstract.html}
\bibitem{lundberg2017}
S. M. Lundberg and S.-I. Lee, ``A unified approach to interpreting model predictions,'' in \emph{Advances in Neural Information Processing Systems 30}, Curran Associates, Inc., 2017.~\href{https://papers.nips.cc/paper/2017/hash/8a20a8621978632d76c43dfd28b67767-Abstract.html}{Official NeurIPS proceedings}.

\def\currentbiburl{https://shap.readthedocs.io/en/latest/generated/shap.LinearExplainer.html}
\bibitem{shaplinear}
SHAP developers, ``shap.LinearExplainer,'' \emph{SHAP documentation}, accessed 21 October 2026.~\url{https://shap.readthedocs.io/en/latest/generated/shap.LinearExplainer.html}.

\def\currentbiburl{https://doi.org/10.1111/j.2517-6161.1972.tb00899.x}
\bibitem{cox1972}
D. R. Cox, ``Regression models and life-tables,'' \emph{Journal of the Royal Statistical Society: Series B (Methodological)}, vol. 34, no. 2, pp. 187--202, 1972, doi:~\url{https://doi.org/10.1111/j.2517-6161.1972.tb00899.x}.

\end{thebibliography}
